\documentclass[11pt]{exam}


\usepackage[exercices,eleve]{profnsi}
\profnsisetup{
	niveau=Terminale,
	matiere=NSI,
	titre=TD — Systèmes d’exploitation et gestion des processus,
	chapitre=11
}

\begin{document}
	\creerDoc
	
	\exercice{Retrouver la bonne commande}
	
	On se place dans un terminal Linux classique.
	
	\begin{enumerate}
		\item Écrire la commande qui permet de :
		\begin{enumerate}
			\item afficher le chemin du répertoire courant ;
			\item lister les fichiers du répertoire courant ;
			\item créer un répertoire nommé \texttt{NSI} ;
			\item supprimer le fichier \texttt{test.txt} ;
			\item afficher le contenu du fichier \texttt{notes.txt}.
		\end{enumerate}
		
		\item On se trouve dans le répertoire \texttt{/home/eleve}.  
		Écrire la commande permettant de :
		\begin{enumerate}
			\item se déplacer dans le répertoire \texttt{/home/eleve/Documents} ;
			\item revenir dans le répertoire parent ;
			\item se déplacer dans le répertoire \texttt{/home/eleve/Téléchargements}.
		\end{enumerate}
	\end{enumerate}
	
	\correction{
		\begin{enumerate}
			\item
			\begin{enumerate}
				\item Afficher le chemin du répertoire courant : \texttt{pwd}
				\item Lister les fichiers du répertoire courant : \texttt{ls} (ou \texttt{ls -l} pour la version détaillée)
				\item Créer un répertoire nommé \texttt{NSI} : \texttt{mkdir NSI}
				\item Supprimer le fichier \texttt{test.txt} : \texttt{rm test.txt}
				\item Afficher le contenu du fichier \texttt{notes.txt} : \texttt{cat notes.txt}
			\end{enumerate}
			
			\item Depuis \texttt{/home/eleve} :
			\begin{enumerate}
				\item Se déplacer dans \texttt{/home/eleve/Documents} :\\
				\texttt{cd /home/eleve/Documents}
				\item Revenir dans le répertoire parent (\texttt{/home}) :\\
				\texttt{cd ..}
				\item Se déplacer dans \texttt{/home/eleve/Téléchargements} :\\
				\texttt{cd /home/eleve/Téléchargements}
			\end{enumerate}
		\end{enumerate}
	}
	
	
	\exercice{Lire un extrait de \texttt{ps}}
	
	On considère l’extrait simplifié suivant d’une commande \texttt{ps -ef} :
	
	\begin{center}
		\begin{tabular}{|c|c|l|}
			\hline
			\textbf{PID} & \textbf{PPID} & \textbf{Commande} \\ \hline
			1 & 0 & /sbin/init \\ \hline
			245 & 1 & /usr/lib/systemd/systemd-logind \\ \hline
			880 & 1 & /usr/sbin/NetworkManager \\ \hline
			1320 & 1 & /usr/bin/gnome-shell \\ \hline
			2045 & 1320 & /usr/bin/firefox \\ \hline
			2101 & 1320 & /usr/bin/gimp \\ \hline
			2300 & 2045 & /usr/lib/firefox/plugin-container \\ \hline
		\end{tabular}
	\end{center}
	
	\begin{enumerate}
		\item Que signifient les colonnes \textbf{PID} et \textbf{PPID} ?
		\item Quel est le \textbf{PID} du processus \texttt{firefox} ?
		\item Quel processus a lancé \texttt{firefox} ? Justifier avec les PID/PPID.
		\item Quels sont les \textbf{fils} du processus \texttt{gnome-shell} ?
		\item Donner l’\textbf{ancêtre commun} de \texttt{firefox} et \texttt{gimp} parmi les processus du tableau.
		\item Représenter sous forme d’\textbf{arbre} les relations père/fils des processus ci-dessus.
	\end{enumerate}
	
	\correction{
		\begin{enumerate}
		\item \textbf{PID} : identifiant du processus (Process ID).  
		\textbf{PPID} : identifiant du processus père (Parent Process ID).
		\item Le PID de \texttt{firefox} est \textbf{2045}.
		\item Le processus qui a lancé \texttt{firefox} est celui dont le PID vaut le PPID de \texttt{firefox}, donc le processus de PID \textbf{1320}, c’est-à-dire \texttt{/usr/bin/gnome-shell}.
		\item Les fils de \texttt{gnome-shell} (PID 1320) sont les processus dont le PPID vaut 1320, donc :
		\begin{itemize}
			\item PID 2045 : \texttt{/usr/bin/firefox}
			\item PID 2101 : \texttt{/usr/bin/gimp}
		\end{itemize}
		\item L’ancêtre commun de \texttt{firefox} et \texttt{gimp} dans ce tableau est \texttt{gnome-shell} (PID \textbf{1320}), qui est leur père direct. (On peut aussi remonter à \texttt{init} PID 1, mais \texttt{gnome-shell} est l’ancêtre commun le plus proche.)
		\item L'arbre :
		\end{enumerate}
		
		— 0\\
		
		\hspace{2cm}— 1 : /sbin/init\\
		
		\hspace{4cm}— 245 : systemd-logind\\
		
		\hspace{4cm}— 880 : NetworkManager\\
		
		\hspace{4cm}— 1320 : gnome-shell\\
		
		\hspace{6cm}— 2045 : firefox\\
		
		\hspace{8cm}— 2300 : plugin-container\\
		
		\hspace{6cm}— 2101 : gimp\\
	}
	
	\exercice{FCFS / SJF}
	
	On considère les processus suivants, avec leur temps d’arrivée et leur durée d’exécution (en unités de temps CPU) :
	
	\begin{center}
		\begin{tabular}{|c|c|c|}
			\hline 
			Processus & Temps d'exécution & Temps d'arrivée \\ 
			\hline 
			A & 5 & 0 \\ 
			\hline 
			B & 2 & 1 \\ 
			\hline 
			C & 3 & 3 \\ 
			\hline 
			D & 4 & 5 \\ 
			\hline 
			E & 4 & 6 \\ 
			\hline 
		\end{tabular} 
	\end{center}
	
	\begin{enumerate}
		\item Compléter le tableau de temps CPU suivant pour représenter l’ordonnancement \textbf{FCFS}
		\item Faites de même pour \textbf{SJF}
		\item Faites de même pour \textbf{le tourniquet} avec des quantum de 1 puis 3
		\item Comparer FCFS et SJF : dans quel cas le temps d’attente moyen vous semble-t-il le plus faible ? Pourquoi ?
	\end{enumerate}
	
	\correction{
		On rappelle les processus :
		
		\begin{center}
			\begin{tabular}{|c|c|c|}
				\hline 
				Processus & Temps d'exécution & Temps d'arrivée \\ 
				\hline 
				A & 5 & 0 \\ 
				\hline 
				B & 2 & 1 \\ 
				\hline 
				C & 3 & 3 \\ 
				\hline 
				D & 4 & 5 \\ 
				\hline 
				E & 4 & 6 \\ 
				\hline 
			\end{tabular} 
		\end{center}
		
		\begin{enumerate}
			\item \textbf{FCFS (First Come, First Served)}\\
			On exécute chaque processus complètement dans l’ordre d’arrivée (en tenant compte des temps d’arrivée) :
			\begin{itemize}
				\item $t=0$ à $t=5$ : A (arrivé à 0)
				\item $t=5$ à $t=7$ : B (arrivé à 1)
				\item $t=7$ à $t=10$ : C (arrivé à 3)
				\item $t=10$ à $t=14$ : D (arrivé à 5)
				\item $t=14$ à $t=18$ : E (arrivé à 6)
			\end{itemize}
			Séquence CPU : \texttt{AAAAA BB CCC DDDD EEEE}
			
			\item \textbf{SJF (Shortest Job First) non préemptif}\\
			À chaque fois que le processeur est libre, on choisit le processus disponible de plus courte durée.
			Avec ces données, l’ordre obtenu est en fait le même que pour FCFS :
			\begin{itemize}
				\item A puis B puis C puis D puis E
			\end{itemize}
			Séquence CPU : \texttt{AAAAA BB CCC DDDD EEEE}
			
			(Dans d’autres jeux de données, SJF et FCFS peuvent produire des ordres très différents.)
			
			\item \textbf{Tourniquet (Round Robin)} avec quantum = 1 puis 3.  
			On suppose un ordonnancement préemptif, en respectant les temps d’arrivée.
			
			\begin{itemize}
				\item \textbf{Quantum = 1} (unité de temps) :
				
				Séquence possible (par unités de temps) :\\
				\texttt{A B A B C A D C E A D C E A D E D E}\\
				(on voit bien l’alternance très fréquente entre processus)
				
				\item \textbf{Quantum = 3} :
				
				Séquence possible :\\
				\texttt{A A A B B C C C A A D D D E E E D E}\\
				(le processeur reste plus longtemps sur le même processus avant de passer au suivant)
			\end{itemize}
			
			\item \textbf{Comparaison FCFS / SJF}\\
			Ici, FCFS et SJF donnent la même séquence, donc le temps d’attente moyen est identique.  
			En général :
			\begin{itemize}
				\item SJF tend à diminuer le temps d’attente moyen,
				\item mais il peut pénaliser les processus longs (famine possible),
				\item alors que FCFS est plus simple mais peut laisser de très longs processus bloquer les autres.
			\end{itemize}
		\end{enumerate}
	}
	
	
	%------------------------------------------------------------
	\exercice{Dîner des philosophes}
	
	On rappelle le principe du problème des \textbf{philosophes} :
	
	\begin{itemize}
		\item cinq philosophes sont assis autour d’une table ;
		\item entre chaque philosophe se trouve une fourchette ;
		\item pour manger, un philosophe a besoin de deux fourchettes (gauche et droite) ;
		\item chaque philosophe alterne entre réfléchir et essayer de manger.
	\end{itemize}
	
	\begin{enumerate}
		\item Expliquer comment une situation d’interblocage peut apparaître dans ce problème.
		\item Proposer une règle simple pour réduire le risque d’interblocage (par exemple en changeant la façon de prendre les fourchettes).
		\item Dans la vraie vie informatique, donner un exemple de situation qui ressemble au problème des philosophes.
	\end{enumerate}
	
	
	\correction{
		\begin{enumerate}
			\item \textbf{Interblocage possible} :  
			Si tous les philosophes prennent d’abord la fourchette de gauche, puis attendent la fourchette de droite, on peut se retrouver dans la situation suivante :
			\begin{itemize}
				\item chaque philosophe détient une fourchette ;
				\item chacun attend la seconde (à droite), qui est déjà tenue par son voisin ;
				\item personne ne peut avancer, tout le monde est bloqué : il y a interblocage.
			\end{itemize}
			
			\item \textbf{Règle simple pour réduire le risque} (plusieurs réponses possibles) :
			\begin{itemize}
				\item imposer à un philosophe de prendre d’abord la fourchette droite, les autres prenant d’abord celle de gauche (on casse la symétrie) ;
				\item limiter le nombre de philosophes autorisés à prendre des fourchettes simultanément (par exemple au plus 4 sur 5) ;
				\item imposer un ordre global sur les fourchettes et les prendre toujours dans le même ordre.
			\end{itemize}
			
			\item \textbf{Exemple réel} :
			\begin{itemize}
				\item Deux programmes qui verrouillent chacun un fichier différent puis essaient de verrouiller le fichier de l’autre ;
				\item Deux transactions de base de données qui verrouillent des tables dans un ordre incompatible ;
				\item Deux processus qui détiennent chacun une ressource matérielle (imprimante, scanner) et attendent l’autre.
			\end{itemize}
		\end{enumerate}
	}
	
	\exercice{Implémentation python}
	
	Écrire le programme python qui gère l'ordonnancement du tourniquet en affichant tour à tour les processus qui ont accès au CPU. Les informations des processus est contenu dans un dictionnaire tel que :
	
	\begin{python}
P = {
	"A" : {
		"t_exec" : 5,
		"t_arrivee" : 0,
	},
	"B" : {
		"t_exec" : 2,
		"t_arrivee" : 1,
	},
	"C" : {
		"t_exec" : 3,
		"t_arrivee" : 3,
	},
	"D" : {
		"t_exec" : 4,
		"t_arrivee" : 5,
	},
	"E" : {
		"t_exec" : 4,
		"t_arrivee" : 6,
	}
}\end{python}



\end{document}
